\section{Further results}\label{app:further}

This appendix holds the per-target detail behind \cref{sec:results,sec:ablations,sec:interp}: every toy target's $\nstar$ for every learner, the post-hoc core included (\cref{app:further-targets}), and what the fitted models contain, target by target (\cref{app:interp-detail}). Apart from the rows of \code{dlheld7}, none of it is held out for DFG v1.3 (\cref{tab:status}).

% ------------------------------------------------------------------------------------------------------------------
\subsection{Every toy target}\label{app:further-targets}

\Cref{tab:targets-a,tab:targets-b,tab:targets-c} give, for each of the 138 targets of the 11 toy suites, $\nstar$ of DFG alone and of every reference, and the post-hoc core's $\nstar$ next to DFG alone's on the same seeds \status{Numerical}. The suite summaries of \cref{tab:suites,tab:pooled,tab:lowdeg,tab:core} are means over these rows, and the targets' definitions are in \repo{study/boolfuncs.py} (for \code{dlheld7}, in \cref{tab:dlheld7}).

\begin{table}[p]
\centering{\scriptsize\setlength{\tabcolsep}{2pt}\renewcommand{\arraystretch}{0.98}
\begin{tabular*}{\linewidth}{@{\extracolsep{\fill}}lrrrrrlrrrr@{}}
\toprule
target & \hd{DFG\\alone} & v3 & \code{lowdeg} & \code{lin} & \code{hgb} & \code{nets_best} (variant) & \hd{DFG alone vs\\\code{nets_best}} & \hd{DFG alone\\vs v3} & \hd{core\\(seeds 0, 1)} & \hd{core vs DFG alone\\(seeds 0, 1)}\\
\midrule
\multicolumn{11}{@{}l}{\code{designed}: 21 targets; development (v1.1, v1.2, v1.3); networks: \code{nets}, \code{nets_wd}}\\
dict & 22 & 16 & 16 & 27 & -- & 40 (wd) & $-0.85$ & $+0.48$ & 16 & $-0.65$\\
and3 & 60 & 25 & 158 & 68 & -- & 70 (wd) & $-0.21$ & $+1.30$ & 37 & $-1.09$\\
maj3 & 56 & 31 & 85 & 56 & -- & 91 (wd) & $-0.70$ & $+0.82$ & 49 & $-0.48$\\
maj9 & 81 & 80 & 323 & 80 & -- & 149 (wd) & $-0.88$ & $+0.02$ & 110 & $+0.39$\\
thr10 & 146 & 121 & 730 & 113 & -- & 307 (wd) & $-1.07$ & $+0.27$ & 154 & $+0.27$\\
par2 & 16 & 16 & 28 & never & -- & 109 (wd) & $-2.77$ & $+0.00$ & 16 & $+0.00$\\
par3 & 28 & 55 & 28 & never & -- & 443 (wd) & $-3.99$ & $-0.99$ & 28 & $+0.01$\\
par4 & 52 & 56 & 28 & never & -- & 1726 (wd) & $-5.06$ & $-0.10$ & 3124 & $+6.02$\\
stair5 & 111 & 109 & 225 & never & -- & 871 (wd) & $-2.97$ & $+0.02$ & 103 & $-0.08$\\
majpar & 61 & 60 & 97 & never & -- & 1026 (wd) & $-4.08$ & $+0.00$ & 49 & $-0.34$\\
parmaj & 111 & 110 & 418 & never & -- & 434 (wd) & $-1.96$ & $+0.02$ & 111 & $-0.00$\\
mux & 63 & 70 & 148 & never & -- & 729 (wd) & $-3.53$ & $-0.14$ & 58 & $-0.28$\\
dnf3 & 178 & 118 & 532 & never & -- & 353 (wd) & $-0.99$ & $+0.60$ & 195 & $-0.02$\\
dlist5 & 61 & 32 & 173 & 85 & -- & 85 (wd) & $-0.49$ & $+0.93$ & 52 & $-0.16$\\
rmaj9 & 167 & 163 & 403 & never & -- & 1273 (wd) & $-2.93$ & $+0.03$ & 209 & $+0.32$\\
andxor & 52 & 40 & 83 & never & -- & 204 (wd) & $-1.98$ & $+0.37$ & 44 & $-0.15$\\
sparsepoly & 102 & 103 & 185 & never & -- & 1238 (wd) & $-3.61$ & $-0.03$ & 82 & $-0.45$\\
rand0 & 25 & 16 & 29 & never & -- & 125 (wd) & $-2.35$ & $+0.62$ & 16 & $-0.73$\\
rand1 & 104 & 95 & 130 & never & -- & 316 (wd) & $-1.60$ & $+0.13$ & 63 & $-0.72$\\
rand2 & 28 & 16 & 16 & 50 & -- & 59 (wd) & $-1.06$ & $+0.83$ & 28 & $+0.13$\\
rand3 & 116 & 114 & 995 & never & -- & 692 (wd) & $-2.58$ & $+0.01$ & 136 & $+0.26$\\
\midrule
\multicolumn{11}{@{}l}{\code{random}: 20 targets; development (v1.1); networks: \code{nets}, \code{nets_wd}}\\
rf00 & 32 & 40 & 16 & 83 & 60 & 62 (wd) & $-0.95$ & $-0.33$ & 30 & $-0.10$\\
rf01 & 52 & 48 & 85 & never & 120 & 256 (wd) & $-2.31$ & $+0.11$ & 74 & $+0.51$\\
rf02 & 51 & 49 & 98 & never & 228 & 191 (wd) & $-1.90$ & $+0.07$ & 49 & $-0.07$\\
rf03 & 91 & 62 & 294 & never & 139 & 276 (wd) & $-1.60$ & $+0.54$ & 108 & $+0.24$\\
rf04 & 54 & 48 & 97 & never & 111 & 203 (wd) & $-1.91$ & $+0.16$ & 56 & $+0.04$\\
rf05 & 75 & 62 & 145 & never & 250 & 493 (wd) & $-2.72$ & $+0.26$ & 72 & $-0.04$\\
rf06 & 31 & 16 & 16 & 37 & 58 & 50 (wd) & $-0.67$ & $+0.97$ & 16 & $-0.97$\\
rf07 & 104 & 88 & 231 & never & 168 & 281 (wd) & $-1.44$ & $+0.23$ & 97 & $-0.10$\\
rf08 & 50 & 21 & 146 & 69 & 53 & 62 (wd) & $-0.32$ & $+1.22$ & 49 & $-0.02$\\
rf09 & 51 & 45 & 37 & never & 292 & 182 (wd) & $-1.83$ & $+0.18$ & 37 & $-0.46$\\
rf10 & 43 & 36 & 86 & 66 & 53 & 79 (wd) & $-0.89$ & $+0.26$ & 37 & $-0.21$\\
rf11 & 28 & 56 & 28 & never & 416 & 445 (wd) & $-4.00$ & $-1.00$ & 28 & $-0.00$\\
rf12 & 97 & 87 & 289 & never & 524 & 610 (wd) & $-2.65$ & $+0.16$ & 87 & $-0.16$\\
rf13 & 54 & 48 & 150 & never & 213 & 245 (wd) & $-2.18$ & $+0.15$ & 59 & $+0.13$\\
rf14 & 53 & 53 & 97 & never & 378 & 404 (wd) & $-2.92$ & $-0.00$ & 55 & $+0.06$\\
rf15 & 51 & 22 & 108 & 67 & 53 & 75 (wd) & $-0.55$ & $+1.23$ & 37 & $-0.48$\\
rf16 & 44 & 36 & 74 & never & 214 & 200 (wd) & $-2.19$ & $+0.27$ & 49 & $+0.16$\\
rf17 & 49 & 56 & 28 & never & never & 1715 (wd) & $-5.14$ & $-0.19$ & 3550 & $+6.19$\\
rf18 & 51 & 46 & 74 & never & 93 & 112 (wd) & $-1.13$ & $+0.16$ & 56 & $+0.12$\\
rf19 & 57 & 59 & 296 & never & 366 & 654 (wd) & $-3.53$ & $-0.06$ & 74 & $+0.39$\\
\bottomrule
\end{tabular*}
}
\caption{Every toy target, part 1 of 3: \code{designed} and \code{random} \status{Numerical}. $\nstar$: the training-set size at which the seed-mean test error reaches 10\% (3 seeds; 2 on \code{random}), interpolated in $\log n$; ``never'': not by $n=16384$; ``--'': not run on the suite. DFG alone, v3 and \code{lowdeg} see the inputs permuted per fit, \code{lin} and \code{hgb} in the natural order. \code{nets_best}: the per-target minimum over the network variants run on the suite (named in each suite's line), with the variant attaining it in brackets: wd, cewd, cewd0.3, cewd1.0 for \code{nets_wd}, \code{nets_cewd}, \code{nets_cewd0.3}, \code{nets_cewd1.0}. Ratios are $\log_2$ of $\nstar$, negative when the first-named learner needs less data; $\pm\infty$: only one of the two reaches 10\% error; n/a: neither. Core (post hoc): seeds 0 and 1, against DFG alone on the same seeds and training sets (\cref{sec:core}). Each suite's line gives its status for DFG v1.3 as in \repo{study/part11/PROTOCOL.md}. Source: \repo{study/part11/RESULTS.json} (\texttt{q2\_suites}, \texttt{q3\_lowdeg}, \texttt{q5b\_core}) via \repo{paper_dfg/make_tables.py} (\texttt{tab\_targets}).}
\label{tab:targets-a}
\end{table}

\begin{table}[p]
\centering{\scriptsize\setlength{\tabcolsep}{2pt}\renewcommand{\arraystretch}{0.98}
\begin{tabular*}{\linewidth}{@{\extracolsep{\fill}}lrrrrrlrrrr@{}}
\toprule
target & \hd{DFG\\alone} & v3 & \code{lowdeg} & \code{lin} & \code{hgb} & \code{nets_best} (variant) & \hd{DFG alone vs\\\code{nets_best}} & \hd{DFG alone\\vs v3} & \hd{core / DFG\\alone (seeds 0, 1)} & \hd{core vs\\DFG alone}\\
\midrule
\multicolumn{11}{@{}l}{\code{dlheld}: 12 targets; development (v1.1); networks: \code{nets}, \code{nets_wd}}\\
maj15 & 130 & 105 & 622 & 110 & 1278 & 210 (wd) & $-0.69$ & $+0.31$ & 206 / 122 & $+0.76$\\
thr30 & 158 & 149 & 1820 & 148 & 1645 & 311 (wd) & $-0.97$ & $+0.09$ & 623 / 158 & $+1.98$\\
majxor & 225 & 224 & 1724 & never & 4076 & 792 (wd) & $-1.82$ & $+0.01$ & 225 / 223 & $+0.02$\\
majmaj5 & 351 & 360 & 5997 & never & 1638 & 2355 (wd) & $-2.75$ & $-0.04$ & 352 / 353 & $-0.00$\\
par5 & 56 & 839 & never & never & never & 7096 (wd) & $-6.99$ & $-3.91$ & never / 56 & $+\infty$\\
mux6 & 197 & 123 & 339 & never & 188 & 358 (wd) & $-0.86$ & $+0.67$ & 210 / 208 & $+0.01$\\
stair6 & 207 & 190 & 1075 & never & 695 & 1495 (wd) & $-2.85$ & $+0.13$ & 205 / 210 & $-0.04$\\
inter2 & 441 & 346 & 2204 & never & 416 & 454 (wd) & $-0.04$ & $+0.35$ & 540 / 518 & $+0.06$\\
majpar5 & 128 & 127 & 198 & never & 4127 & 4648 (wd) & $-5.19$ & $+0.01$ & 108 / 165 & $-0.62$\\
hrf0 & 80 & 62 & 32 & 191 & 63 & 133 (wd) & $-0.74$ & $+0.36$ & 46 / 77 & $-0.75$\\
hrf1 & 119 & 116 & 121 & never & 461 & 789 (wd) & $-2.73$ & $+0.04$ & 123 / 116 & $+0.09$\\
hrf2 & 46 & 54 & 26 & 93 & 60 & 90 (wd) & $-0.98$ & $-0.23$ & 56 / 50 & $+0.16$\\
\midrule
\multicolumn{11}{@{}l}{\code{dlheld2}: 7 targets; held out for an earlier version only; networks: \code{nets}, \code{nets_wd}}\\
thr20 & 138 & 124 & 1054 & 117 & 1262 & 304 (wd) & $-1.14$ & $+0.16$ & 455 / 134 & $+1.76$\\
maj21 & 112 & 109 & 693 & 116 & 1956 & 216 (wd) & $-0.94$ & $+0.04$ & 174 / 109 & $+0.67$\\
xmaj & 61 & 66 & 168 & never & 1538 & 871 (wd) & $-3.83$ & $-0.10$ & 56 / 62 & $-0.15$\\
dnf4x2 & 235 & 112 & 467 & never & 98 & 306 (wd) & $-0.38$ & $+1.07$ & 252 / 228 & $+0.14$\\
sym8 & 54 & 53 & 1553 & never & 2689 & 377 (wd) & $-2.81$ & $+0.03$ & 100 / 54 & $+0.88$\\
h2rf0 & 119 & 103 & 118 & never & 220 & 389 (wd) & $-1.71$ & $+0.21$ & 112 / 117 & $-0.06$\\
h2rf1 & 82 & 31 & 179 & 103 & 58 & 103 (wd) & $-0.33$ & $+1.40$ & 80 / 79 & $+0.00$\\
\midrule
\multicolumn{11}{@{}l}{\code{dlheld3}: 10 targets; held out for an earlier version only; networks: \code{nets}, \code{nets_cewd}, \code{nets_wd}}\\
cnf3 & 189 & 126 & 600 & never & 164 & 361 (cewd) & $-0.93$ & $+0.58$ & 164 / 196 & $-0.26$\\
sym9 & 111 & 111 & 3267 & never & 7845 & 350 (cewd) & $-1.65$ & $+0.00$ & 111 / 111 & $-0.00$\\
dlist8 & 39 & 24 & 29 & 91 & 61 & 90 (wd) & $-1.21$ & $+0.69$ & 30 / 55 & $-0.86$\\
thrint & 186 & 139 & 659 & 137 & 243 & 190 (cewd) & $-0.03$ & $+0.42$ & 287 / 260 & $+0.14$\\
par3or & 71 & 63 & 177 & never & 2009 & 725 (wd) & $-3.35$ & $+0.18$ & 60 / 69 & $-0.21$\\
poly4 & 176 & 177 & 195 & never & never & 2745 (cewd) & $-3.96$ & $-0.01$ & 3127 / 161 & $+4.28$\\
mixmaj & 114 & 96 & 289 & never & 368 & 326 (cewd) & $-1.52$ & $+0.24$ & 89 / 103 & $-0.22$\\
mux11 & 509 & 466 & 1486 & never & 573 & 1245 (wd) & $-1.29$ & $+0.13$ & 850 / 637 & $+0.42$\\
h3rf0 & 108 & 46 & 16 & 112 & 60 & 96 (wd) & $+0.17$ & $+1.22$ & 30 / 120 & $-1.98$\\
h3rf1 & 133 & 59 & 31 & 101 & 63 & 101 (wd) & $+0.39$ & $+1.17$ & 31 / 135 & $-2.13$\\
\midrule
\multicolumn{11}{@{}l}{\code{dlsymdev}: 8 targets; development (v1.2, v1.3); networks: \code{nets}, \code{nets_cewd}, \code{nets_wd}}\\
sym7 & 56 & 55 & 1356 & never & -- & 329 (cewd) & $-2.56$ & $+0.02$ & 55 / 56 & $-0.02$\\
ex3of8 & 123 & 118 & 6509 & never & -- & 664 (cewd) & $-2.43$ & $+0.06$ & 94 / 103 & $-0.14$\\
band11 & 828 & 829 & never & never & -- & 1836 (cewd) & $-1.15$ & $-0.00$ & 889 / 775 & $+0.20$\\
far9 & 98 & 71 & 1433 & never & -- & 296 (wd) & $-1.60$ & $+0.47$ & 86 / 95 & $-0.13$\\
ssym9 & 111 & 104 & 2992 & never & -- & 324 (cewd) & $-1.55$ & $+0.10$ & 111 / 111 & $-0.00$\\
symxpar & 112 & 111 & 806 & never & -- & 1410 (cewd) & $-3.66$ & $+0.00$ & 112 / 112 & $-0.00$\\
twosym & 415 & 415 & 5326 & never & -- & 1817 (cewd) & $-2.13$ & $-0.00$ & 449 / 388 & $+0.21$\\
thrsum & 79 & 80 & 406 & never & -- & 330 (cewd) & $-2.07$ & $-0.02$ & 140 / 112 & $+0.32$\\
\midrule
\multicolumn{11}{@{}l}{\code{dlheld4}: 12 targets; exposed (v1.2 built after exposure); networks: \code{nets}, \code{nets_cewd}, \code{nets_wd}}\\
sym11 & 167 & 160 & 7778 & never & never & 374 (cewd) & $-1.16$ & $+0.06$ & 187 / 193 & $-0.04$\\
ex4of9 & 205 & 194 & 6580 & never & never & 411 (cewd) & $-1.01$ & $+0.08$ & 193 / 202 & $-0.07$\\
tail13 & 134 & 100 & 3992 & never & never & 420 (cewd) & $-1.65$ & $+0.42$ & 148 / 157 & $-0.09$\\
band10 & 440 & 432 & never & never & never & 3440 (cewd) & $-2.97$ & $+0.02$ & 432 / 440 & $-0.02$\\
ssym8 & 148 & 177 & 3873 & never & 8072 & 396 (cewd) & $-1.42$ & $-0.27$ & 100 / 105 & $-0.07$\\
symxor & 406 & 405 & 2569 & never & 6909 & 946 (wd) & $-1.22$ & $+0.00$ & 461 / 439 & $+0.07$\\
symand & 149 & 151 & 761 & never & 1307 & 370 (cewd) & $-1.31$ & $-0.02$ & 100 / 126 & $-0.33$\\
sumpar & 181 & 164 & 316 & never & never & 1837 (cewd) & $-3.34$ & $+0.14$ & 146 / 164 & $-0.17$\\
majsym & 102 & 90 & 978 & never & 2840 & 595 (cewd) & $-2.54$ & $+0.17$ & 96 / 98 & $-0.03$\\
dnf3x3 & 193 & 115 & 538 & never & 148 & 373 (wd) & $-0.95$ & $+0.74$ & 212 / 191 & $+0.15$\\
h4rf0 & 55 & 47 & 98 & never & 160 & 206 (wd) & $-1.90$ & $+0.22$ & 55 / 55 & $+0.01$\\
h4rf1 & 96 & 54 & 114 & 158 & 61 & 143 (wd) & $-0.58$ & $+0.82$ & 29 / 103 & $-1.84$\\
\bottomrule
\end{tabular*}
}
\caption{Every toy target, part 2 of 3: \code{dlheld}, \code{dlheld2}, \code{dlheld3}, \code{dlsymdev} and \code{dlheld4} \status{Numerical}. Columns as in \cref{tab:targets-a}.}
\label{tab:targets-b}
\end{table}

\begin{table}[p]
\centering{\scriptsize\setlength{\tabcolsep}{2pt}\renewcommand{\arraystretch}{0.98}
\begin{tabular*}{\linewidth}{@{\extracolsep{\fill}}lrrrrrlrrrr@{}}
\toprule
target & \hd{DFG\\alone} & v3 & \code{lowdeg} & \code{lin} & \code{hgb} & \code{nets_best} (variant) & \hd{DFG alone vs\\\code{nets_best}} & \hd{DFG alone\\vs v3} & \hd{core\\(seeds 0, 1)} & \hd{core vs DFG alone\\(seeds 0, 1)}\\
\midrule
\multicolumn{11}{@{}l}{\code{dlheld5}: 12 targets; exposed (v1.2 built after exposure); networks: \code{nets}, \code{nets_cewd}, \code{nets_wd}}\\
h5s0 & 201 & 205 & never & never & never & 803 (wd) & $-2.00$ & $-0.03$ & 216 & $+0.11$\\
h5s1 & never & never & never & never & never & never & n/a & n/a & never & n/a\\
h5s2 & 598 & 1154 & 6456 & never & 10860 & 1850 (cewd) & $-1.63$ & $-0.95$ & 678 & $+0.88$\\
h5s3 & 1642 & 1757 & never & never & never & 6454 (cewd) & $-1.97$ & $-0.10$ & 1752 & $-0.01$\\
h5s4 & 222 & 184 & 1766 & never & 1465 & 1801 (wd) & $-3.02$ & $+0.27$ & 178 & $-0.06$\\
h5s5 & 420 & 421 & never & never & never & 12254 (cewd) & $-4.87$ & $-0.00$ & 449 & $+0.10$\\
h5s6 & 597 & 573 & never & never & never & 3963 (cewd) & $-2.73$ & $+0.06$ & 354 & $-0.07$\\
h5s7 & 6186 & 6760 & never & never & never & 7641 (cewd) & $-0.30$ & $-0.13$ & 1698 & $-0.01$\\
h5t0 & 109 & 105 & 402 & never & 2886 & 352 (cewd) & $-1.68$ & $+0.06$ & 185 & $+0.72$\\
h5t1 & 100 & 95 & 223 & never & 993 & 325 (cewd) & $-1.70$ & $+0.07$ & 100 & $-0.06$\\
h5rf0 & 215 & 203 & 318 & never & 2658 & 1547 (cewd) & $-2.85$ & $+0.08$ & 164 & $-0.84$\\
h5rf1 & 98 & 97 & 476 & never & 9281 & 2512 (cewd) & $-4.68$ & $+0.01$ & 149 & $+0.60$\\
\midrule
\multicolumn{11}{@{}l}{\code{dlheld6}: 12 targets; development (v1.3); networks: \code{nets}, \code{nets_cewd}, \code{nets_cewd0.3}, \code{nets_cewd1.0}, \code{nets_wd}}\\
h6\_00 & never & never & never & never & never & 12304 (cewd1.0) & $+\infty$ & n/a & never & n/a\\
h6\_01 & 431 & 430 & 7470 & never & never & 1655 (cewd1.0) & $-1.94$ & $+0.00$ & 428 & $-0.01$\\
h6\_02 & 384 & 389 & 666 & never & never & 996 (cewd0.3) & $-1.37$ & $-0.02$ & 414 & $-0.01$\\
h6\_03 & 129 & 98 & 270 & never & 103 & 231 (cewd1.0) & $-0.84$ & $+0.40$ & 188 & $+0.59$\\
h6\_04 & 79 & 29 & 259 & 96 & 59 & 96 (wd) & $-0.27$ & $+1.44$ & 65 & $-0.29$\\
h6\_05 & 42 & 33 & 16 & 48 & 59 & 56 (wd) & $-0.39$ & $+0.38$ & 27 & $-0.55$\\
h6\_06 & 347 & 349 & never & never & never & 1872 (cewd0.3) & $-2.43$ & $-0.01$ & 387 & $-0.03$\\
h6\_07 & 448 & 447 & 5684 & never & 4492 & 1812 (cewd1.0) & $-2.02$ & $+0.00$ & 495 & $+0.35$\\
h6\_08 & 479 & 449 & 2542 & never & 5940 & 1720 (cewd0.3) & $-1.84$ & $+0.09$ & 522 & $-0.16$\\
h6\_09 & 391 & 380 & 7850 & never & 2706 & 744 (cewd0.3) & $-0.93$ & $+0.04$ & 227 & $-0.69$\\
h6\_10 & 213 & 206 & never & never & never & 451 (cewd1.0) & $-1.08$ & $+0.05$ & 649 & $+1.60$\\
h6\_11 & 192 & 184 & 2993 & never & 2264 & 357 (cewd1.0) & $-0.90$ & $+0.06$ & 156 & $-0.14$\\
\midrule
\multicolumn{11}{@{}l}{\code{dlfamdev}: 12 targets; development (v1.3); networks: none run}\\
fd\_00 & 138 & 75 & 239 & -- & -- & -- & -- & $+0.87$ & 139 & $+0.29$\\
fd\_01 & 82 & 79 & 189 & -- & -- & -- & -- & $+0.05$ & 86 & $+0.09$\\
fd\_02 & 61 & 66 & 31 & -- & -- & -- & -- & $-0.11$ & 31 & $-0.93$\\
fd\_03 & 120 & 81 & 205 & -- & -- & -- & -- & $+0.56$ & 116 & $-0.04$\\
fd\_04 & 170 & 171 & never & -- & -- & -- & -- & $-0.00$ & 111 & $+0.00$\\
fd\_05 & 158 & 147 & 2555 & -- & -- & -- & -- & $+0.11$ & 167 & $+0.11$\\
fd\_06 & 355 & 318 & 9543 & -- & -- & -- & -- & $+0.16$ & 223 & $-0.46$\\
fd\_07 & 134 & 98 & 344 & -- & -- & -- & -- & $+0.45$ & 105 & $-0.23$\\
fd\_08 & 51 & 29 & 89 & -- & -- & -- & -- & $+0.81$ & 53 & $+0.04$\\
fd\_09 & 104 & 101 & 3210 & -- & -- & -- & -- & $+0.04$ & 99 & $-0.08$\\
fd\_10 & never & never & never & -- & -- & -- & -- & n/a & never & n/a\\
fd\_11 & 114 & 123 & 61 & -- & -- & -- & -- & $-0.10$ & 61 & $-0.84$\\
\midrule
\multicolumn{11}{@{}l}{\code{dlheld7}: 12 targets; held out (the only suite); networks: \code{nets_cewd0.3}, \code{nets_cewd1.0}, \code{nets_wd}}\\
f7\_00 & 211 & 203 & 7061 & never & 6208 & 1545 (cewd0.3) & $-2.87$ & $+0.06$ & 211 & $-0.01$\\
f7\_01 & 119 & 106 & 595 & 111 & 684 & 119 (cewd1.0) & $-0.00$ & $+0.17$ & 195 & $+0.32$\\
f7\_02 & 190 & 106 & 435 & never & 120 & 337 (wd) & $-0.83$ & $+0.85$ & 186 & $-0.01$\\
f7\_03 & 352 & 352 & 4205 & never & 2199 & 1599 (cewd1.0) & $-2.18$ & $-0.00$ & 344 & $-0.04$\\
f7\_04 & 97 & 70 & 121 & never & 84 & 164 (wd) & $-0.76$ & $+0.47$ & 99 & $+0.01$\\
f7\_05 & never & never & never & never & never & never & n/a & n/a & never & n/a\\
f7\_06 & 82 & 44 & 179 & 127 & 59 & 108 (wd) & $-0.39$ & $+0.90$ & 82 & $+0.02$\\
f7\_07 & 167 & 174 & 1123 & never & 1373 & 415 (cewd1.0) & $-1.32$ & $-0.06$ & 121 & $-0.44$\\
f7\_08 & 644 & 634 & never & never & never & 2932 (cewd1.0) & $-2.19$ & $+0.02$ & 454 & $+0.03$\\
f7\_09 & 222 & 215 & 591 & never & never & 883 (cewd0.3) & $-1.99$ & $+0.05$ & 355 & $+0.68$\\
f7\_10 & 84 & 82 & 336 & 87 & 114 & 119 (cewd1.0) & $-0.50$ & $+0.03$ & 104 & $+0.23$\\
f7\_11 & 51 & 55 & 29 & 119 & 63 & 114 (wd) & $-1.16$ & $-0.12$ & 30 & $-1.07$\\
\bottomrule
\end{tabular*}
}
\caption{Every toy target, part 3 of 3: \code{dlheld5}, \code{dlheld6}, \code{dlfamdev} and the held-out \code{dlheld7} \status{Numerical}. Columns as in \cref{tab:targets-a}. \code{dlfamdev} has no network, \code{lin} or \code{hgb} runs.}
\label{tab:targets-c}
\end{table}

% ------------------------------------------------------------------------------------------------------------------
% Appended for sections/interp.tex (part 1 of "What it learns"); tables from make_interp.py.
\subsection{What the models contain: the detail}\label{app:interp-detail}

This subsection holds the per-target detail behind \cref{sec:interp-atlas,sec:interp-calib}: every atlas target (\cref{tab:interp-atlas-targets}), the spurious terms by $n$ (\cref{tab:interp-spur}), the identity checks of the analyses, DFG's choice of variant, the kernel background and the cells where a model loses structure as $n$ grows. All of it is seed 0, from \repo{study/part11/interp/ATLAS.md} and \repo{study/part11/interp/calib/CALIB.md} with their verifications.

\begin{table}[tbp]
\centering{\scriptsize\setlength{\tabcolsep}{3pt}\renewcommand{\arraystretch}{0.97}
\begin{tabular}{@{}llrrrrrrcrrc@{}}
\toprule
& & & & \multicolumn{5}{c}{first $n$ with test error $\le0.1$} & \multicolumn{3}{c}{$n=16384$}\\
\cmidrule(lr){5-9}\cmidrule(lr){10-12}
target & suite & $n^*_0$ ($n^*_3$) & complete & $n$ & error & var. & size & T / R / S & error & size & T / R / S\\
\midrule
\multicolumn{12}{@{}l}{\emph{Parities}}\\
par2 & designed & 16 (16) & 16 & 16 & 0 & f & 2 & 1.00 / 0.00 / 0.00 & 0 & 2 & 1.00 / 0.00 / 0.00\\
par4 & designed & 56 (52) & 64 & 64 & 0 & d & 2 & 1.00 / 0.00 / 0.00 & 0 & 2 & 1.00 / 0.00 / 0.00\\
par5 & dlheld & 56 (56) & 64 & 64 & 0 & d & 2 & 1.00 / 0.00 / 0.00 & 0 & 2 & 1.00 / 0.00 / 0.00\\
\addlinespace[2pt]
\multicolumn{12}{@{}l}{\emph{Staircases}}\\
stair5 & designed & 108 (111) & 128 (512) & 128 & 0 & f & 7 & 0.99 / 0.00 / 0.00 & 0 & 17 & 1.00 / 0.00 / 0.00\\
stair6 & dlheld & 210 (207) & 256 (--) & 256 & 0 & d & 8 & 0.97 / 0.00 / 0.00 & 0 & 81 & 0.96 / 0.21 / 0.00\\
\addlinespace[2pt]
\multicolumn{12}{@{}l}{\emph{Majorities of parities; andxor}}\\
majpar & designed & 57 (61) & 64 & 64 & 0.067 & f & 4 & 0.92 / 0.00 / 0.00 & 0 & 11 & 0.92 / 0.22 / 0.00\\
majpar5 & dlheld & 182 (128) & 256 & 256 & 0.011 & d & 6 & 0.98 / 0.00 / 0.00 & 0 & 29 & 0.99 / 0.11 / 0.00\\
andxor & designed & 63 (52) & 256 & 64 & 0.092 & f & 3 & 0.85 / 0.00 / 0.00 & 0 & 4 & 1.00 / 0.00 / 0.00\\
\addlinespace[2pt]
\multicolumn{12}{@{}l}{\emph{Multiplexers}}\\
mux6 & dlheld & 213 (197) & 1024 & 256 & 0.002 & f & 43 & 0.89 / 0.63 / 0.16 & 0 & 17 & 1.00 / 0.00 / 0.00\\
mux11 & dlheld3 & 430 (509) & 4096 & 512 & 0 & d & 76 & 0.83 / 0.56 / 0.02 & 0 & 65 & 1.00 / 0.00 / 0.00\\
\addlinespace[2pt]
\multicolumn{12}{@{}l}{\emph{DNF, CNF, decision lists, read-once formulas}}\\
cnf3 & dlheld3 & 187 (189) & 1024 & 256 & 0.004 & f & 48 & 0.85 / 0.48 / 0.09 & 0 & 67 & 0.92 / 0.10 / 0.07\\
f7\_02 & dlheld7 & 188 (190) & 2048 & 256 & 0 & f & 43 & 0.91 / 0.30 / 0.12 & 0 & 69 & 0.98 / 0.04 / 0.00\\
dlist8 & dlheld3 & 71 (39) & 64 (--) & 128 & 0.049 & d & 40 & 0.89 / 0.02 / 0.36 & 0 & 149 & 0.91 / 0.04 / 0.10\\
f7\_04 & dlheld7 & 99 (97) & 4096 & 128 & 0.021 & d & 33 & 0.96 / 0.00 / 0.24 & 0 & 181 & 0.92 / 0.09 / 0.05\\
f7\_06 & dlheld7 & 80 (82) & 2048 & 128 & 0.024 & d & 42 & 0.87 / 0.50 / 0.28 & 0 & 84 & 0.93 / 0.10 / 0.09\\
\addlinespace[2pt]
\multicolumn{12}{@{}l}{\emph{Random formulas}}\\
rf06 & random & 46 (31) & -- & 64 & 0.045 & d & 2 & 0.96 / 0.00 / 0.00 & 0 & 66 & 0.98 / 0.00 / 0.03\\
hrf0 & dlheld & 75 (80) & 16384 & 128 & 0 & f & 8 & 0.97 / 0.37 / 0.00 & 0 & 32 & 1.00 / 0.00 / 0.00\\
\addlinespace[2pt]
\multicolumn{12}{@{}l}{\emph{Majorities and halfspaces}}\\
maj9 & designed & 85 (81) & 64 & 128 & 0 & d & 37 & 1.00 / 0.00 / 0.01 & 0 & 54 & 1.00 / 0.00 / 0.00\\
thrint & dlheld3 & 163 (186) & 64 (1024) & 256 & 0.071 & d & 34 & 1.08 / 0.39 / 0.22 & 0 & 43 & 0.85 / 0.28 / 0.00\\
thr30 & dlheld & 187 (158) & 64 & 256 & 0.058 & d & 31 & 1.00 / 0.00 / 0.00 & 0.009 & 149 & 0.91 / 0.15 / 0.00\\
f7\_01 & dlheld7 & 177 (119) & 64 & 256 & 0 & d & 56 & 0.75 / 0.00 / 0.27 & 0 & 57 & 0.92 / 0.30 / 0.00\\
\addlinespace[2pt]
\multicolumn{12}{@{}l}{\emph{Parities of two majorities}}\\
parmaj & designed & 111 (111) & 128 & 128 & 0 & f & 37 & 1.00 / 0.00 / 0.00 & 0 & 19 & 1.00 / 0.00 / 0.00\\
majxor & dlheld & 223 (225) & 256 & 256 & 0.0005 & d & 9 & 0.96 / 0.00 / 0.10 & 0 & 137 & 1.00 / 0.00 / 0.00\\
\addlinespace[2pt]
\multicolumn{12}{@{}l}{\emph{Majority of majorities}}\\
majmaj5 & dlheld & 348 (351) & -- & 512 & 0.023 & d & 49 & 0.91 / 0.69 / 0.11 & 0 & 256 & 0.90 / 0.20 / 0.00\\
\addlinespace[2pt]
\multicolumn{12}{@{}l}{\emph{Symmetric and signed sums}}\\
sym9 & dlheld3 & 111 (111) & 128 & 128 & 0 & d & 4 & 1.00 / 0.00 / 0.00 & 0 & 75 & 0.88 / 0.00 / 0.21\\
band11 & dlsymdev & 891 (828) & 1024 & 1024 & 0 & d & 5 & 1.00 / 0.00 / 0.00 & 0 & 62 & 0.98 / 0.12 / 0.00\\
f7\_08 & dlheld7 & 448 (644) & 512 & 512 & 0.012 & f & 52 & 0.96 / 0.00 / 0.20 & 0 & 60 & 1.05 / 0.18 / 0.01\\
h5s7 & dlheld5 & 1705 (6186) & 2048 (8192) & 2048 & 0.0007 & d & 164 & 0.79 / 0.43 / 0.07 & 0.0005 & 255 & 0.59 / 0.45 / 0.05\\
\addlinespace[2pt]
\multicolumn{12}{@{}l}{\emph{Signed-sum mixtures}}\\
f7\_07 & dlheld7 & 205 (167) & 256 (--) & 256 & 0.006 & d & 21 & 1.00 / 0.00 / 0.00 & 0 & 87 & 1.00 / 0.00 / 0.00\\
f7\_09 & dlheld7 & 222 (222) & 256 (--) & 256 & 0.0005 & f & 15 & 1.00 / 0.00 / 0.00 & 0 & 35 & 1.00 / 0.00 / 0.00\\
symand & dlheld4 & 173 (149) & -- & 256 & 0.012 & d & 36 & 1.00 / 0.00 / 0.19 & 0 & 80 & 1.00 / 0.00 / 0.00\\
f7\_10 & dlheld7 & 90 (84) & 256 & 128 & 0 & d & 33 & 1.18 / 0.64 / 0.12 & 0 & 27 & 1.00 / 0.00 / 0.00\\
f7\_00 & dlheld7 & 221 (211) & 256 (--) & 256 & 0 & f & 21 & 1.00 / 0.00 / 0.00 & 0 & 101 & 0.98 / 0.04 / 0.00\\
f7\_03 & dlheld7 & 352 (352) & 512 (--) & 512 & 0.008 & d & 68 & 0.85 / 0.45 / 0.01 & 0 & 110 & 1.00 / 0.00 / 0.00\\
f7\_11 & dlheld7 & 72 (51) & -- & 128 & 0.016 & d & 8 & 0.97 / 0.00 / 0.17 & 0 & 115 & 0.97 / 0.11 / 0.06\\
\addlinespace[2pt]
\multicolumn{12}{@{}l}{\emph{Residue classes mod 3}}\\
f7\_05 & dlheld7 & never (never) & -- & \multicolumn{5}{c}{never} & 0.3245 & 1 & 0.00 / 0.00 / 0.00\\
h6\_00 & dlheld6 & never (never) & -- & \multicolumn{5}{c}{never} & 0.3430 & 1 & 0.00 / 0.00 / 0.00\\
\bottomrule
\end{tabular}
}
\caption{Every target of the structure atlas: DFG alone, seed 0 \status{Numerical}. $n^*_0$ ($n^*_3$): the 10\% test-error point of the seed-0 fits (of the seed mean in the result file; 2 seeds for rf06), interpolated in $\log n$. ``Complete'': the first grid $n$ at which some compact representation of the target is complete in the model and, in brackets where it differs, the $n$ from which it stays complete (--: it does not stay complete to $n=16384$). Variant: d = default, f = forced, as DFG's cross-validation chose it. Size: features, the constant included. T / R / S: the shares of target, relevant non-target and spurious terms in the model's output, each the standard deviation of the terms' summed contribution over that of the output; parts that partly cancel can exceed 1. Families as in \cref{tab:interp-atlas}. Source: \repo{study/part11/interp/tables/atlas_family.json} and \texttt{atlas\_targets.csv}, via \repo{paper_dfg/make_interp.py}; the leading terms of every model are in \repo{study/part11/interp/tables/atlas_summary.md}.}
\label{tab:interp-atlas-targets}
\end{table}

\begin{table}[tbp]
\centering{\scriptsize\setlength{\tabcolsep}{3.6pt}
\begin{tabular}{@{}rrrrrrrrr@{}}
\toprule
& & & \multicolumn{3}{c}{terms reading an irrelevant input} & \multicolumn{2}{c}{spurious share of the output} & test error\\
\cmidrule(lr){4-6}\cmidrule(lr){7-8}
$n$ & \hd{fits with\\error $\le0.1$} & \hd{admitted\\terms} & all (share) & \hd{linear-block\\characters} & \hd{others (share of\\the other terms)} & mean & max & \hd{with / without\\them (mean)}\\
\midrule
16 & 2 & 5 & 1 (0.200) & 0 & 1/5 (0.200) & 0.024 & 1.000 & 0.3713 / 0.3654\\
32 & 7 & 18 & 3 (0.167) & 0 & 3/18 (0.167) & 0.033 & 0.752 & 0.2818 / 0.2800\\
64 & 22 & 362 & 250 (0.691) & 247 & 3/68 (0.044) & 0.136 & 0.863 & 0.1158 / 0.0958\\
128 & 38 & 399 & 204 (0.511) & 185 & 19/164 (0.116) & 0.063 & 0.493 & 0.0190 / 0.0191\\
256 & 41 & 405 & 174 (0.430) & 159 & 15/212 (0.071) & 0.025 & 0.170 & 0.0027 / 0.0028\\
512 & 41 & 403 & 103 (0.256) & 92 & 11/300 (0.037) & 0.016 & 0.111 & 0.0007 / 0.0006\\
1024 & 41 & 408 & 96 (0.235) & 76 & 20/323 (0.062) & 0.012 & 0.148 & 0.0002 / 0.0002\\
2048 & 41 & 519 & 53 (0.102) & 50 & 3/464 (0.006) & 0.003 & 0.086 & 0.0000 / 0.0000\\
4096 & 41 & 681 & 56 (0.082) & 50 & 6/625 (0.010) & 0.006 & 0.067 & 0.0000 / 0.0000\\
8192 & 41 & 608 & 52 (0.086) & 50 & 2/553 (0.004) & 0.003 & 0.055 & 0.0000 / 0.0000\\
16384 & 41 & 691 & 55 (0.080) & 52 & 3/632 (0.005) & 0.003 & 0.060 & 0.0000 / 0.0000\\
\midrule
all & & 4499 & 1047 (0.233) & 961 & 86/3364 (0.026) & & & \\
\bottomrule
\end{tabular}
}
\caption{Spurious terms on the 21 \code{designed} and 20 \code{random} targets, seed 0, 41 fits per $n$ \status{Numerical}. A term is spurious if it reads an input on which the target does not depend; a sum block counts as one term. ``Linear-block characters'': irrelevant degree-1 characters that entered with the linear block. Spurious share: the standard deviation of the spurious terms' summed contribution over that of the model's output. Test error with and without the spurious terms' contributions, mean over the 41 fits; from $n=128$ on the two differ by at most 0.0002 (\repo{study/part11/interp/calib/VERIFY_CALIB.md}). Source: \repo{study/part11/interp/calib/calib.json} (\texttt{A2.by\_n}), via \repo{paper_dfg/make_interp.py}.}
\label{tab:interp-spur}
\end{table}

\paragraph{The analyses' fits are the evaluation's} \status{Numerical}. Every analysis refits DFG through the evaluation's own code path and checks the refits against the result files, which store test errors and FLOP counts but no models.
\begin{itemize}
\item \emph{The atlas}: 517 cells at $d=30$ (47 targets of all 11 suites, every $n$, seed 0) and 66 cells at $d=15$, 60 and 120, equal in test error and FLOP count (\repo{study/part11/interp/tables/identity.md}).
\item \emph{The order analysis}: 363 fits (11 \code{designed} targets, every $n$, seeds 0--2), equal in both; the residual rebuilt from the recorded weights matches DFG's own round log in all of them (\repo{study/part11/interp/order/ORDER.md}).
\item \emph{The theory checks}: 330 new seed-0 cells of \code{designed} and \code{random}, equal in both, and the 110 fits on null labels, whose round-0 decisions an independent replica reproduces in all 110 (\repo{study/part11/interp/theory/THEORY.md}).
\item \emph{The false-admission and compute analysis} runs an instrumented copy of the frozen file with 64 inserted lines that only read. Its verifier rebuilt the copy byte for byte and found that removing the inserted statements leaves a syntax tree identical to the frozen file's. Every fit checked against the frozen module or the result files is equal (88 on noise, 451 on the designed and random targets, 35 profiled cells), and the unchecked fits are equal by that construction (\repo{study/part11/interp/calib/VERIFY_CALIB.md}).
\item \emph{This paper's check} of h5s1 (\cref{sec:interp-fail}): 6 refits, equal in test error (\repo{paper_dfg/check_h5s1.json}).
\end{itemize}

\paragraph{DFG's own choice of variant} \status{Numerical}. At $n\le512$ the paired cross-validation (\cref{sec:smalln}) kept the default variant in 204 of the atlas's 222 cells and chose the forced one in 18; in 12 of these 18 the forced model is the target's first fit (par2 at 16, andxor and majpar at 64, stair5, parmaj and hrf0 at 128, and six more at 256--512). Refitting the variant DFG did not choose on the same rows (\repo{study/part11/interp/cv_counterfactual.py}), the two differ in test error in 140 cells, and the choice made is the better one in 89 of them. The forced variant would have been better in 48 of the 204 default cells, most of all par4 at $n=32$ (0 against 0.488): its forced character is right on all 32 rows, but the 10-fold fits on 28--29 rows do not search degree 4. The mean test error over the 222 cells is 0.2257 for DFG's choice, 0.2319 for the default always, 0.2256 for the forced variant always and 0.2111 for the better of the two in each cell. \emph{Reading} \status{Heuristic}. The paired test is conservative where the folds are too small for the forced guess to be right, which is where the guess helps most. There v3's stacker, which kept both variants, did better (\cref{sec:interp-atlas}).

\paragraph{Below the structure} \status{Numerical}. In the 198 cells (36 targets) below the first $n$ at which a compact representation is complete, 187 models are the default variant and 11 the forced one. 108 are the constant alone, with the background; 87 hold at least one true character; the median test error is 0.31 (\repo{study/part11/interp/tables/atlas_below.csv}).

\paragraph{The kernel background} \status{Numerical}. The background is added after the weights are fitted, so removing it from a model gives exactly DFG's prediction without it. Removed, the test error is lower in 85 of the 407 cells, higher in 54 and unchanged in 268. It helps most at $n=16$--32, where the model holds little else and the target has a strong degree-1 part (maj9, thrint, thr30, majmaj5, symand, f7\_04, f7\_10, rf06): without it the error would be 0.19--0.27 higher (maj9 at $n=32$: 0.491 against 0.283). It hurts where the residual's degree-1 correlations are noise: on the residue classes at $n=16$--32 (h6\_00 at 16: 0.488 with it, 0.343 without), and at the first fits of andxor and majpar at $n=64$ (0.092 and 0.067 against 0).

\paragraph{Structure lost as $n$ grows} \status{Numerical}. On one training set per $n$ the models are not monotone in $n$. A compact representation present at one $n$ is missing at a larger one, with the test error at or near 0, in stair5 (its degree-5 prefix character at $n=256$ only), stair6 (its degree-6 prefix character from 512 on), thrint (at 512), dlist8 (at 256 and 16384), f7\_07 and f7\_09 (complete only at 256), f7\_03 (only at 512) and f7\_00 (from 4096 on). h5s7 loses its whole block at 4096 (error 0.368) and recovers it at 8192. After the first fit the test error rises with $n$ in nine targets, by at most 0.013 (f7\_02: 0 at $n=256$, 0.013 at 512). In the order analysis, DFG drops 1139 of the 2157 character-seeds it carried at a smaller $n$, because each fit stops once every training sign is right (\cref{sec:interp-order}).
